\documentclass[11pt]{article}

%% Code-level companion to advstd_techniques.tex.
%%
%% This guide focuses on:
%%   (1) the top-ranked safe combination zono + prev_pgd + tau=0.5;
%%   (2) the hyper-attack (PGD warm-start) integration in advstd Phase 2;
%%   (3) the VHAGaR dependency constraints in advstd Phase 2.
%%
%% All paths are relative to /root/Downloads/for_dana/vaghar_org/.
%% Line numbers reflect the source as of 2026-04-25.

\usepackage[margin=1in]{geometry}
\usepackage{amsmath, amssymb}
\usepackage{enumitem}
\usepackage{xcolor}
\usepackage{booktabs}
\usepackage{listings}
\usepackage{hyperref}
\usepackage{courier}

\hypersetup{
  colorlinks=true,
  linkcolor=blue!50!black,
  urlcolor=blue!50!black,
  citecolor=blue!50!black,
}

\lstdefinestyle{jl}{
  basicstyle=\ttfamily\footnotesize,
  breaklines=true,
  frame=single,
  framesep=4pt,
  numbers=left,
  numberstyle=\tiny\color{gray},
  keywordstyle=\color{blue!60!black},
  commentstyle=\color{green!40!black},
  stringstyle=\color{red!60!black},
  showstringspaces=false,
  columns=fullflexible,
  language=Java, % closest to Julia syntax for listings
}

\newcommand{\src}[1]{\texttt{\small #1}}
\newcommand{\tech}[1]{\textsf{#1}}
\newcommand{\flag}[1]{\texttt{-{}-#1}}
\newcommand{\fileline}[2]{\src{#1}:\src{#2}}

\title{Implementation Guide: \\
  the safe combo \texttt{zono} + \texttt{prev\_pgd} + $\tau=0.5$, \\
  with hyper-attack and VHAGaR-deps wiring}
\author{Code-level companion to \texttt{advstd\_techniques.tex}}
\date{2026-04-25}

\begin{document}
\maketitle

\begin{abstract}
This document is a code-level walkthrough of three things: (1)~the
top-ranked safe combination \texttt{zono} + \texttt{prev\_pgd} +
$\tau=0.5$ from Section~9 of \texttt{advstd\_techniques.tex}; (2)~the
hyper-attack (PGD warm-start) integration into advstd Phase~2;
(3)~the VHAGaR dependency constraints (\flag{activate\_vaghgar\_deps})
as they apply to advstd Phase~2. The aim is operational: every claim is
backed by a file path and line range, and the call order in
\src{main\_advanced\_standard} (\src{run.jl}) is reproduced verbatim so
a reader can step through it.

The combo is sound: at solver \texttt{OPTIMAL} it returns
$\delta_{\text{exact}}$. The other two pieces (hyper-attack and
VHAGaR-deps) are independent toggles that work alongside the combo.
Hyper-attack provides a Cutoff and (in \texttt{prev\_pgd} mode) a
warm-start that the consensus filter then rationalises against
$N_1$. VHAGaR-deps adds inter-copy constraints linking the org and
pert encodings and is preserved exactly as in standard mode.
\end{abstract}

\tableofcontents

%%─────────────────────────────────────────────────────────────────────
\section{Code Map}
\label{sec:codemap}

All paths relative to \src{vaghar\_org/}.

\begin{center}
\footnotesize
\begin{tabular}{@{}l l l@{}}
\toprule
Topic & File & Key entry points (line) \\
\midrule
Phase-2 driver           & \src{run.jl}                                        & \src{main\_advanced\_standard} (630--1052) \\
Phase-2 N2-only wrapper  & \src{run.jl}                                        & \src{main\_advanced\_standard\_n2} (1181--1188) \\[3pt]
Tech 1 zono Source A     & \src{utils/perturbation\_intervals.jl}              & \src{compute\_diff\_bounds\_zonotope} (1317--1740) \\
Tech 1 zono Source B     & \src{utils/perturbation\_intervals.jl}              & \src{compute\_n2\_zono\_abs\_bounds} (1743--1988) \\
Tech 1 bound integration & \src{utils/MIPVerify.jl/.../core\_ops.jl}            & \src{intersect\_per\_copy\_bounds} (199--246) \\
Tech 2 prev\_pgd hints   & \src{utils/help\_functions.jl}                      & \src{apply\_n1\_var\_hints!} (598--741) \\
Tech 2 \emph{p}-formula  & \src{utils/help\_functions.jl}                      & \src{compute\_varhint} (507--523) \\
Tech 3 $\tau$ decision   & \src{utils/n2\_relax\_decision.jl}                  & \src{compute\_n2\_relax\_decision!} (23--130) \\
Tech 3 MIP emit          & \src{utils/MIPVerify.jl/.../core\_ops.jl}            & \src{relu} (248--610), gating at 554--605 \\[3pt]
Hyper-attack subprocess  & \src{utils/hyper\_attack.jl}                        & \src{hyper\_attack} (98--113) \\
Hyper-attack PGD code    & \src{utils/hyper\_attack.py}                        & \src{attack} (120--303) \\
Hyper-attack hint apply  & \src{utils/hyper\_attack.jl}                        & \src{hyper\_attack\_hints} (1--35) \\
Cutoff plumbing          & \src{utils/mip.jl}                                  & \src{mip\_set\_attr} (20--34) \\[3pt]
VHAGaR deps (entry)      & \src{utils/perturbation\_dependencies.jl}           & \src{perturbation\_dependencies} (202--232) \\
VHAGaR deps (encoder)    & \src{utils/perturbation\_dependencies.jl}           & \src{encode\_dependencies} (112--200) \\
\bottomrule
\end{tabular}
\end{center}

%%─────────────────────────────────────────────────────────────────────
\section{Phase-2 call sequence (the source of truth)}
\label{sec:driver}

The advstd Phase-2 driver \src{main\_advanced\_standard} in
\fileline{run.jl}{630--1052} processes one $(c_s, c_t)$ pair at a time.
Globally (before the pair loop), it loads the \texttt{n1\_state}
directory and calls \src{compute\_diff\_bounds\_zonotope} +
\src{compute\_n2\_zono\_abs\_bounds} once.\footnote{Grep
\src{run.jl} around lines 1090--1100 for the once-only zonotope
computation. The Source-B pass is genuinely shared across $(c_s, c_t)$
because it depends only on $N_2$'s weights and the input box.}
Per pair, the order is:

\begin{enumerate}[leftmargin=2em, itemsep=2pt, label=\arabic*.]
  \item Run \src{hyper\_attack} for $N_2$ to get
    \src{suboptimal\_solution\_n2} for the Gurobi
    \texttt{Cutoff}.~\fileline{run.jl}{909--912}
  \item \src{mip\_reset()}; install $N_1$'s neuron bounds globally
    via \src{set\_n1\_neuron\_bounds(n1\_layers\_info)}.~\fileline{run.jl}{930--935}
  \item \emph{(Tech 3, only if $\tau\ge 0$)} call
    \src{compute\_n2\_relax\_decision!($\tau$)}
    --- this populates the \src{n2\_relax\_decision} dict consulted
    later inside \src{relu()}.~\fileline{run.jl}{945--949}
  \item \src{get\_model(...)} builds the dual-copy $N_2$ MIP. Inside
    every \src{relu()} call, both the per-copy bound integration
    (Sources A/B/C) and the triangle-vs-bigM choice happen
    here.~\fileline{run.jl}{951--956}
  \item Build the result filename suffix from active flags.
    For our combo: \src{\_BoundTightPertRelax0.5\_zonoBounds\_varHintPrevPGD}.
    \fileline{run.jl}{962--989}
  \item \emph{(hyper-attack)} \src{hyper\_attack\_hints(m\_n2,
    token*"\_n2", c\_tag, c\_target)} — read PGD output from
    \src{/tmp} and call \src{set\_start\_value} on the matched binaries.
    Filename suffix \src{\_HyperAttackHints} appended.~\fileline{run.jl}{999--1006}
  \item \emph{(VHAGaR deps)} \src{perturbation\_dependencies(m\_n2,
    nn2, ...)}; suffix \src{\_VagharDeps}.~\fileline{run.jl}{1007--1011}
  \item \emph{(perturbed intervals)}
    \src{perturbed\_interval\_constraints(m\_n2, nn2, "org",
    "perturbation")}; suffix \src{\_PerturbedIntervals}.~\fileline{run.jl}{1012--1015}
  \item \src{mip\_set\_delta\_property}, \src{set\_optimizer},
    \src{mip\_set\_attr} (this is where the \texttt{Cutoff} is
    installed using the value from step~1).~\fileline{run.jl}{1016--1018}
  \item \emph{(Tech 2)} call
    \src{apply\_n1\_var\_hints!(m\_n2, var\_hint\_mode, ...)} ---
    must run \emph{after} \src{hyper\_attack\_hints} so the
    Start-consensus filter sees PGD's values.~\fileline{run.jl}{1027--1034}
  \item \src{optimize!(m\_n2)}, then \src{mip\_log}.~\fileline{run.jl}{1043--1044}
\end{enumerate}

\noindent\textbf{The order matters.} Two specific dependencies:

\begin{itemize}[noitemsep, leftmargin=2em]
  \item \src{compute\_n2\_relax\_decision!} must run \emph{before}
    \src{get\_model}, because \src{relu()} consults the global
    decision dict at MIP-build time. Calling it after the build is a
    no-op.
  \item \src{apply\_n1\_var\_hints!} (the \texttt{prev\_pgd} branch)
    must run \emph{after} \src{hyper\_attack\_hints}, otherwise PGD
    silently overwrites the consensus filter's writes.
\end{itemize}

%%─────────────────────────────────────────────────────────────────────
\section{Tech 1 --- \texttt{zono} bound tightening}
\label{sec:zono}

\subsection{Source A --- diff zonotope $N_1 \to N_2$}

\paragraph{What it produces.}
For every ReLU neuron $(m, k)$ in $N_2$, two arrays:
\src{relu\_diff\_up\_bounds[m][k]} and
\src{relu\_diff\_down\_bounds[m][k]}, sound over-approximations of
$d^{N_2}_{m,k} - d^{N_1}_{m,k}$. These are stored as 1-indexed Julia
arrays of length $K$ = number of ReLU layers per copy
(\src{help\_functions.jl}, declared near line 74).

\paragraph{Where.}
\src{compute\_diff\_bounds\_zonotope} in
\fileline{utils/perturbation\_intervals.jl}{1317--1740}. The affine
update at lines 1397--1423 implements
\[
  c_{m,k} \leftarrow \sum_{k'} w^2_{m,k,k'}\, c_{m-1,k'}
    + \sum_{k'} \Delta w_{m,k,k'}\, \mathrm{mid}_{k'} + \Delta b_{m,k},
\]
\[
  G'_{\mathrm{corr}} = W_2 \cdot G_{m-1}, \qquad
  G'_{\mathrm{indep}} = \Delta W \cdot \mathrm{diag}(\mathrm{rad}_{1..K'}),
\]
\[
  G_m \leftarrow [G'_{\mathrm{corr}} \;\big|\; G'_{\mathrm{indep}}],
\]
with $\mathrm{mid}_{k'} = \tfrac12(\max(0, u^{N_1}_{m{-}1,k'}) +
\max(0, l^{N_1}_{m{-}1,k'}))$ and $\mathrm{rad}_{k'}$ the
half-width (post-ReLU clipping is applied, matching
\texttt{advstd\_techniques.tex} \S 4.1.2). The DeepZ ReLU update
(formulas $\lambda_k = u_k/(u_k - l_k)$, $\mu_k = -u_k l_k /
(2(u_k - l_k))$) lives at \src{perturbation\_intervals.jl} lines
1614--1615; the per-split-neuron new column with $\mu_k$ at row $k$
and zero elsewhere is appended at line 1958.

\paragraph{Caching.}
Source A is computed exactly once per Phase-2 invocation, before the
$(c_s, c_t)$ loop. The resulting global arrays are then reused for
every pair.

\subsection{Source B --- absolute $N_2$ zonotope}

\paragraph{What it produces.}
\src{n2\_abs\_up\_bounds[m][k]} and \src{n2\_abs\_down\_bounds[m][k]}
--- sound over-approximations of $\hat z^{N_2}_{m,k}$ propagated
through $N_2$ alone, with a single zonotope whose initial hull
covers \emph{both} the clean input box $[0,1]^n$ and the perturbed
input box $[-\varepsilon, 1{+}\varepsilon]^n$. Initialisation:
$\mathbf c^{\mathrm{B}}_0 = \tfrac12\mathbf 1$,
$G^{\mathrm{B}}_0 = \mathrm{diag}(\tfrac12 + \varepsilon)$
(\fileline{utils/perturbation\_intervals.jl}{1839--1843}). The
intersection with Source A is applied \emph{at every ReLU layer},
\emph{before} the DeepZ relaxation (line 1921), so the relaxed tube
narrows as Source A tightens it.

\subsection{Per-copy bound integration --- \src{intersect\_per\_copy\_bounds}}

\paragraph{Where.}
\fileline{utils/MIPVerify.jl/src/net\_components/core\_ops.jl}{199--246}.
Called by \emph{both} the decision site (\src{compute\_n2\_relax\_decision!})
and the emit site (\src{relu()}).

\begin{lstlisting}[style=jl]
function intersect_per_copy_bounds(
    l_init::Real, u_init::Real,
    nn_layer::Int, nn_neuron::Int,
    m_idx::Int, k_idx::Int,
    version::AbstractString,
)::Tuple{Float64, Float64}
    l = Float64(l_init); u = Float64(u_init)
    # Source A: N1 per-copy bound + diff bound
    if !isempty(n1_neuron_bounds) && (version == "org" || version == "perturbation")
        key = (nn_layer, nn_neuron)
        if haskey(n1_neuron_bounds, key)
            (n1_u, n1_l) = n1_neuron_bounds[key]
            if !isempty(relu_diff_up_bounds) && 1 <= m_idx <= length(relu_diff_up_bounds) &&
               1 <= k_idx <= length(relu_diff_up_bounds[m_idx])
                u = min(u, n1_u + relu_diff_up_bounds[m_idx][k_idx])
                l = max(l, n1_l + relu_diff_down_bounds[m_idx][k_idx])
            end
        end
    end
    # Source B: shared absolute N2 hull
    if !isempty(n2_abs_up_bounds) && (version == "org" || version == "perturbation")
        if 1 <= m_idx <= length(n2_abs_up_bounds) &&
           1 <= k_idx <= length(n2_abs_up_bounds[m_idx])
            u = min(u, n2_abs_up_bounds[m_idx][k_idx])
            l = max(l, n2_abs_down_bounds[m_idx][k_idx])
        end
    end
    # ... Source C (n1 probe LP) intentionally omitted for our combo
    return (l, u)
end
\end{lstlisting}

The per-copy story: \src{n1\_neuron\_bounds} is keyed by
\src{(neurons\_names.layer, neurons\_names.neuron)}, where
\src{neurons\_names.layer} runs $1..K$ during the org pass and
$K{+}1..2K$ during the pert pass. The diff and Source-B arrays are
indexed by \src{(m\_idx, k\_idx)} with \src{m\_idx} resetting to
$1..K$ in each pass, so they are shared between copies. This
asymmetry is exactly what \texttt{advstd\_techniques.tex} \S 4.3
calls out: only Source A varies per copy.

\paragraph{Emit-site call.}
Inside \src{relu()} at \fileline{core\_ops.jl}{295--301}:

\begin{lstlisting}[style=jl]
if network_version == "org" || network_version == "perturbation"
    (l, u) = intersect_per_copy_bounds(
        l, u,
        neurons_names.layer, neurons_names.neuron,
        layer_counter, neurons_names.neuron,
        network_version,
    )
end
\end{lstlisting}

\src{layer\_counter} is reset to 0 at the start of each pass
(see e.g. \fileline{utils/perturbation\_models.jl}{142}) and
incremented by 1 per ReLU layer at \fileline{core\_ops.jl}{739}, so
\src{layer\_counter} runs $1..K$ for each pass --- matching the
decision-site call argument-for-argument.

%%─────────────────────────────────────────────────────────────────────
\section{Tech 2 --- \texttt{prev\_pgd} variable hints}
\label{sec:prevpgd}

\paragraph{Where.}
\src{apply\_n1\_var\_hints!} in
\fileline{utils/help\_functions.jl}{598--741}; the \texttt{prev}
math in \src{compute\_varhint} at lines 507--523.

\paragraph{What \texttt{prev\_pgd} delivers.}
For each surviving $N_2$ binary $a_i$ (i.e.\ those with
$l^{N_2}_i < 0 < u^{N_2}_i$ after Tech 1's elimination), compute
\(
  \mathtt{hint\_val}(a_i) = (p_i \ge 0.5)\,?\, v^{N_1}_i \,:\, 1 - v^{N_1}_i
\)
using the prev rule, then route it through \src{set\_start\_value}
filtered by PGD's three-way consensus. \emph{No}
\texttt{VarHintVal}/\texttt{VarHintPri} is written.

\subsection{Step 1 --- name-match $N_1$'s primal}

\fileline{help\_functions.jl}{619--654}. The N2 binary's name is
\src{\{nv\}a\_layerCount\{lc\}\_neuronCount\{nc\}\_\{layer\}\_\{neuron\}};
the regex \src{r"a\_layerCount\textbackslash d+\_neuronCount\textbackslash d+\_(\textbackslash d+)\_(\textbackslash d+)\$"}
extracts \src{(layer, neuron)} which keys directly into
\src{layers\_info\_dict} (the global $N_2$ tightened-bounds table)
and into the saved \src{n1\_layers\_info}. Confidence and other
non-ReLU binaries are skipped because their names don't match the
regex.

\subsection{Step 2 --- the prev rule for $p_i$}

\fileline{help\_functions.jl}{671--699}. Two proxies:

\begin{itemize}[noitemsep]
  \item \emph{Active} ($v^{N_1}_i = 1$): read $\hat z^{N_1}_i$
    directly from N1's primal via the matching \src{x\_rect}
    name --- same name pattern, with the literal substring
    \src{a\_layerCount} replaced by \src{x\_rect\_layerCount} once
    (line 681).
  \item \emph{Inactive} ($v^{N_1}_i = 0$): the big-$M$ encoding
    only constrains $\hat z^{N_1}_i \in [l^{N_1}_i, 0]$, so use
    the midpoint $l^{N_1}_i / 2$ as a proxy (line 684).
\end{itemize}

The diff bound is looked up at line 688 via
\src{r = (layer <= K) ? layer : (layer - K)}, mapping the
1..2K global layer back to the 1..K diff slot (the same diff
applies to both copies). Then \src{compute\_varhint} at lines
507--523 builds
\(
  \tilde I_i = [\hat z^{N_1}_i + d^\downarrow_i,\; \hat z^{N_1}_i + d^\uparrow_i] \cap [l^{N_2}_i, u^{N_2}_i]
\)
and computes
\(
  p_i = \mathrm{length}(\tilde I_i \cap \text{N1-side}) / \mathrm{length}(\tilde I_i)
\),
where N1-side is $(-\infty, 0]$ if $v^{N_1}_i = 0$ else
$[0, +\infty)$. Degenerate (empty/point) intervals fall back to
$(\mathtt{hint\_val} = v^{N_1}_i,\, \mathtt{pri} = 1)$ at line 513.

\subsection{Step 3 --- PGD-consensus Start filter}

\fileline{help\_functions.jl}{702--720}:

\begin{lstlisting}[style=jl]
if mode == VH_DIRECT_PGD || mode == VH_PREV_PGD
    v_pgd_raw = JuMP.start_value(v)
    if v_pgd_raw === nothing
        JuMP.set_start_value(v, Float64(hint_val))      # fill
        n_pgd_silent_filled += 1
    else
        v_pgd_bit = Int(round(v_pgd_raw))
        if v_pgd_bit == hint_val
            n_pgd_agreed_nop += 1                       # leave
        else
            JuMP.set_start_value(v, nothing)            # withdraw
            n_pgd_disagreed_withdrew += 1
        end
    end
    applied += 1
else
    MOI.set(JuMP.backend(m_n2), Gurobi.VariableAttribute("VarHintVal"), JuMP.index(v), Float64(hint_val))
    MOI.set(JuMP.backend(m_n2), Gurobi.VariableAttribute("VarHintPri"), JuMP.index(v), hint_pri)
    applied += 1
end
\end{lstlisting}

The non-PGD modes (\texttt{prev}/\texttt{direct}) take the
\texttt{else} branch and write the soft hint pair. The \texttt{\_pgd}
modes never enter that branch, so they never write
\texttt{VarHintVal}/\texttt{VarHintPri} --- the hint priority is
discarded by design.

\paragraph{What survives the filter.}
The function tracks four counters logged at line 728:
\src{pgd-silent-filled}, \src{pgd-agreed-nop},
\src{pgd-disagreed-withdrew}, and \src{flipped} (how many of the
\src{hint\_val}s disagreed with $v^{N_1}_i$ after the prev rule's
flip). The third counter --- withdrawals --- represents binaries
where \emph{both} signals had an opinion but disagreed. Sweep logs
typically show the silent-filled count is the largest, because PGD
only sets Start on a small subset of binaries.

\paragraph{Survivor filter.}
Lines 661--664 short-circuit any binary whose tightened bound is
single-signed:
\begin{lstlisting}[style=jl]
if l_n2 >= 0 || u_n2 <= 0
    n_tier1_skipped += 1
    continue
end
\end{lstlisting}
This is defensive: \src{relu()} would not have created the binary
in the first place. The counter logged at line 728 names this
``Tech2-eliminated-but-survived''.

\paragraph{Abort guard.}
Line 736: if \src{applied == 0} but the model has any binaries, the
function errors out before \src{optimize!} so we never silently ship
a no-hint result mislabeled as \texttt{prev\_pgd}.

%%─────────────────────────────────────────────────────────────────────
\section{Tech 3 --- $\tau=0.5$ BoundTightPertRelax}
\label{sec:tau}

\subsection{Decision site}

\fileline{utils/n2\_relax\_decision.jl}{23--130}. Pseudocode-faithful
to \texttt{advstd\_techniques.tex} \S 4.3. The tiered rule:

\begin{lstlisting}[style=jl]
if max(g_org, g_pert) <= threshold              # Tier 1
    relax_org, relax_pert = true, true
elseif min(g_org, g_pert) <= threshold          # Tier 2
    if g_org <= g_pert
        relax_org, relax_pert = true, false     # org wins on tie
    else
        relax_org, relax_pert = false, true
    end
else                                            # Tier 3
    relax_org, relax_pert = false, false
end
\end{lstlisting}

The gap formula $g(l, u) = u \cdot |l| / (2(u-l))$ is in
\src{\_tri\_gap} at \fileline{core\_ops.jl}{187--188} (returns 0 for
single-signed intervals). The decision starts with
\src{(l\_init, u\_init) = (-Inf, +Inf)} (\src{n2\_relax\_decision.jl}
lines 48--49) and lets \src{intersect\_per\_copy\_bounds} narrow.
This is intentionally conservative: a wider starting envelope can
only shrink the relaxation count.

\paragraph{Decision dict layout.}
Lines 113--123:
\begin{lstlisting}[style=jl]
if relax_org
    n2_relax_decision[(nn_layer_org, k_idx)]  = (true, false)
end
if relax_pert
    n2_relax_decision[(nn_layer_pert, k_idx)] = (false, true)
end
\end{lstlisting}
The keys distinguish copies (\src{nn\_layer\_org = m\_idx},
\src{nn\_layer\_pert = m\_idx + K}). The lookup site (next subsection)
selects the right boolean by \src{network\_version}.

\subsection{MIP emit site}

\fileline{core\_ops.jl}{554--605}. After
\src{intersect\_per\_copy\_bounds} narrowed $(l, u)$:

\begin{lstlisting}[style=jl]
if adv_std_n2_relax_threshold >= 0.0 &&
   (network_version == "org" || network_version == "perturbation") &&
   haskey(n2_relax_decision, (neurons_names.layer, neurons_names.neuron))
    (relax_org, relax_pert) = n2_relax_decision[(neurons_names.layer, neurons_names.neuron)]
    should_relax = (network_version == "org") ? relax_org : relax_pert
    if should_relax
        # emit triangle inequalities on (l, u): the SAME (l, u) the decision saw
        # ... triangle big-M-free encoding here ...
    end
end
\end{lstlisting}

The load-bearing invariant is that the emit site uses the same $(l,
u)$ that the decision evaluated. Both sites compute $(l, u)$ via
\src{intersect\_per\_copy\_bounds} with arg-equivalent inputs (see
\S\ref{sec:zono}). The triangle is therefore emitted on the exact
interval whose gap was tested against $\tau$.

\subsection{Soundness}

For $(m, k, \mathrm{copy}) \notin R$: the exact big-$M$ encoding is
emitted, identical to baseline. For $(m, k, \mathrm{copy}) \in R$:
the triangle inequalities are emitted on
$[l^{N_2,\mathrm{copy}}_{m,k}, u^{N_2,\mathrm{copy}}_{m,k}]$, which
covers the exact ReLU on the same interval. So
$\mathcal F_{2} \subseteq \mathcal F_{2}^{(R)}$ and
$\delta_{\mathrm{BTPR}} \ge \delta_{\mathrm{exact}}$. (Same argument
as \texttt{advstd\_techniques.tex} \S 4.3 \emph{Soundness}.)

%%─────────────────────────────────────────────────────────────────────
\section{Hyper-attack (PGD) in advstd Phase 2}
\label{sec:hyperattack}

\paragraph{Two channels.} \flag{use\_hyper\_attack}\;\texttt{true}
arms two distinct things in Phase 2:

\begin{enumerate}[noitemsep]
  \item A \texttt{Cutoff} on $N_2$'s MIP (always installed when PGD
    succeeded).
  \item Either MIP-Start hints (\texttt{adv\_std\_mip\_start=true},
    rejected by every safe combo) or Start values feeding the
    \texttt{prev\_pgd}/\texttt{direct\_pgd} consensus filter.
\end{enumerate}

\subsection{Subprocess invocation}

\src{hyper\_attack} at \fileline{utils/hyper\_attack.jl}{98--113}
shells out to the Python script:

\begin{lstlisting}[style=jl]
function hyper_attack(dataset, c_tag, c_target, token_signature,
                      model_name, model_path, perturbation, perturbation_size;
                      force_cpu::Bool=false)
    pre_time = @elapsed begin
        cmd_args = ["python3", "./utils/hyper_attack.py",
                    "--dataset", string(dataset),
                    "--source", string(c_tag-1),       # 0-indexed for Python
                    "--target", string(c_target-1),
                    "--token", token_signature,
                    "--model", model_name,
                    "--model_path", model_path*"th",   # .pth checkpoint
                    "--perturbation", perturbation,
                    "--perturbation_size", create_perturbation_string(perturbation_size)]
        if force_cpu; push!(cmd_args, "--cpu"); end
        run(Cmd(cmd_args))
        best_feasible_via_optimization = read_best_val_via_optimization(c_tag, c_target, token_signature)
    end
    return best_feasible_via_optimization, pre_time
end
\end{lstlisting}

The Python side (\src{utils/hyper\_attack.py}) runs PGD for 500
iterations (\src{attack} at lines 120--303), keeps the top-$K$
adversarial inputs satisfying source$\to$target misclassification,
and writes three \src{/tmp} files per $(c_s, c_t)$ pair (or, on
failure, a single \src{fail\_*} file):

\begin{itemize}[noitemsep]
  \item \src{/tmp/booleans\_\{c-1\}\_\{t-1\}\_\{token\}.txt} ---
    comma-separated 0/1/-1 flags per binary (-1 = abstain).
  \item \src{/tmp/strings\_\{c-1\}\_\{t-1\}\_\{token\}.txt} ---
    matching variable names. \src{build\_str} at lines 100--117
    constructs the names: e.g.\
    \src{org a\_layerCount2\_neuronCount0\_2\_5}
    for ReLU layer 2 / neuron 5 in the org copy, or with
    \src{perturbation} for the perturbed copy.
  \item \src{/tmp/best\_val\_\{c-1\}\_\{t-1\}\_\{token\}.txt} ---
    the feasible objective value, used as the Cutoff.
\end{itemize}

In Phase 2, \emph{both} $N_1$ and $N_2$ get their own PGD run with
distinct tokens (\src{token\_signature} for $N_1$ at
\fileline{run.jl}{860--862}, \src{token\_signature * "\_n2"} for
$N_2$ at \fileline{run.jl}{909--912}).

\subsection{Hint application}

\src{hyper\_attack\_hints} at \fileline{utils/hyper\_attack.jl}{1--35}
reads the two text files, matches names against
\src{JuMP.all\_variables(m)}, and calls \src{set\_start\_value} on
matched binaries (skipping \src{-1} abstains and unmatched names):

\begin{lstlisting}[style=jl]
function hyper_attack_hints(m, token, c_tag, c_target)
    av = JuMP.all_variables(m)
    if isfile("/tmp/fail_"*string(c_tag-1)*"_"*string(c_target-1)*"_"*token*".txt")
        rm("/tmp/fail_"*string(c_tag-1)*"_"*string(c_target-1)*"_"*token*".txt")
    else
        # ... read booleans + strings ...
        for n in eachindex(arr_strings)
            if indexes_[n] == -1 || arr_booleans[n] == -1; continue; end
            set_start_value(av[indexes_[n]], arr_booleans[n])
        end
    end
end
\end{lstlisting}

\paragraph{Cutoff plumbing.}
Whatever PGD returned as the best feasible objective rides in
\src{d\_n2[:suboptimal\_solution]}. \src{mip\_set\_attr} at
\fileline{utils/mip.jl}{20--34}:
\begin{lstlisting}[style=jl]
function mip_set_attr(m, perturbation, d, timout)
    ...
    if d[:suboptimal_solution] != 0
        set_optimizer_attribute(m, "Cutoff", d[:suboptimal_solution])
    end
    ...
end
\end{lstlisting}
\texttt{Cutoff} prunes any branch whose objective is provably worse,
giving Gurobi a strong dual-bound prune even when no MIP-Start is
installed.

\subsection{Interaction with \texttt{prev\_pgd}}

The \texttt{prev\_pgd} consensus filter at
\fileline{help\_functions.jl}{702--720} reads
\src{JuMP.start\_value(v)} on every surviving binary --- which is
exactly what \src{hyper\_attack\_hints} just wrote. The filter
reconciles the two signals:

\begin{itemize}[noitemsep]
  \item Binaries where PGD wrote a value and the prev-rule
    \src{hint\_val} agrees: PGD's value is left in place
    (\emph{leave}).
  \item Binaries where they disagree: \src{set\_start\_value(v,
    nothing)} unsets PGD's value (\emph{withdraw}).
  \item Binaries where PGD wrote nothing: filled with the prev-rule
    \src{hint\_val} (\emph{fill}).
\end{itemize}

So in our combo, hyper-attack contributes \emph{(a)} the Cutoff
unconditionally and \emph{(b)} a partial Start that the consensus
filter trims down to the high-confidence overlap with $N_1$'s
transferred bits. If PGD failed entirely (\src{fail\_*} file
present, every \src{start\_value} is \src{nothing}), the combo
degrades gracefully to ``prev fills every gap'' --- i.e.\ a pure
$N_1$-transfer Start.

\paragraph{Order constraint, again.} \src{hyper\_attack\_hints} runs
at \fileline{run.jl}{999--1006}; \src{apply\_n1\_var\_hints!} runs at
\fileline{run.jl}{1027--1034}. Reversing the order would cause PGD to
overwrite the consensus filter's writes and silently break the
combo's semantics.

%%─────────────────────────────────────────────────────────────────────
\section{VHAGaR dependencies in advstd Phase 2}
\label{sec:vagharddeps}

\paragraph{Toggle.} \flag{activate\_vaghgar\_deps}\;\texttt{true}.
The advstd Phase-2 driver wires it into both passes:
\fileline{run.jl}{880--882} for the $N_1$ MIP and
\fileline{run.jl}{1007--1011} for the $N_2$ MIP.

\paragraph{Entry point.}
\src{perturbation\_dependencies} in
\fileline{utils/perturbation\_dependencies.jl}{202--232}:

\begin{lstlisting}[style=jl]
function perturbation_dependencies(m, nn, perturbation, perturbation_size, w, h, k;
                                   activation_start=1, layers_offset=nothing,
                                   perturbation_var=nothing)
    layers_n = isnothing(layers_offset) ? layers_number(nn) : layers_offset
    phi_dep = fill(NaN, (1, w, h, k))
    perturbation_init_deps(phi_dep, perturbation, perturbation_size)
    activation_cnt = activation_start
    if reuse_bounds_conf.is_reuse_bounds_and_deps && activation_start == 1
        for (activation_cnt, phi_dep) in enumerate(reuse_bounds_conf.reusable_deps)
            encode_dependencies(m, layers_n, phi_dep, activation_cnt)
        end
    else
        for l in nn.layers
            if occursin("Flatten", string(typeof(l)))
                phi_dep = phi_dep |> l
            elseif occursin("Linear", string(typeof(l))) || occursin("Conv", string(typeof(l)))
                phi_dep_l = dep_propagation(l, phi_dep)
                phi_dep = dep_additional(m, layers_n, l, phi_dep, phi_dep_l,
                                          perturbation, perturbation_size,
                                          activation_cnt, activation_start, perturbation_var)
            elseif occursin("ReLU", string(typeof(l)))
                phi_dep = encode_dependencies(m, layers_n, phi_dep, activation_cnt)
                if activation_start == 1; push!(reuse_bounds_conf.reusable_deps, phi_dep); end
                if all(isnan, phi_dep); break; end
                activation_cnt += 1
            end
        end
    end
end
\end{lstlisting}

\paragraph{What \src{phi\_dep} encodes.}
A 4D matrix shaped like the perturbation input, holding one of
$\{0,\,1,\,-1,\,\mathrm{NaN}\}$ per element:
\begin{itemize}[noitemsep]
  \item \texttt{0} --- org and pert copies are forced equal at this
    neuron.
  \item \texttt{1} --- pert dominates org (org $\le$ pert).
  \item \texttt{-1} --- org dominates pert.
  \item \texttt{NaN} --- relation unknown; the encoder resolves it via a
    sub-LP if the upstream interval bounds permit (see below).
\end{itemize}

\paragraph{Per-neuron constraint emit.}
\src{encode\_dependencies} at \fileline{utils/perturbation\_dependencies.jl}{112--200}
walks each ReLU layer index \src{activation\_cnt}. For each neuron
$n$ where both copies have a stored bound entry
(\src{layers\_info\_dict[(activation\_cnt, n)]} for org and
\src{layers\_info\_dict[(activation\_cnt + layers\_n, n)]} for pert,
where \src{layers\_n = K} is the per-copy layer count):

\begin{lstlisting}[style=jl]
if dep == 0
    @constraint(m, av[ind_o+1] == av[ind_p+1])    # x_rect equality
    if has_a_o && has_a_p
        @constraint(m, av[ind_o+2] == av[ind_p+2])  # binary equality
    end
elseif dep == 1
    @constraint(m, av[ind_o+1] <= av[ind_p+1])
    if has_a_o && has_a_p
        @constraint(m, av[ind_o+2] <= av[ind_p+2])
    end
elseif dep == -1
    @constraint(m, av[ind_o+1] >= av[ind_p+1])
    if has_a_o && has_a_p
        @constraint(m, av[ind_o+2] >= av[ind_p+2])
    end
end
\end{lstlisting}

The \src{ind\_o + 1} and \src{ind\_o + 2} offsets correspond to the
\src{x\_rect} (post-ReLU) and binary $a$ that follow the
pre-activation \src{ind\_o} in the variable layout. The
\src{has\_a\_o} / \src{has\_a\_p} guards (lines 124--134) check the
expected variable name suffix \src{\_\{layer\}\_\{neuron\}} before
emitting the binary-equality constraint --- this defends against
neurons whose binary was \emph{eliminated} by Tech 1
(\src{l \(\ge\) 0} or \src{u \(\le\) 0}) or \emph{relaxed} by Tech 3
(triangle, no binary). When the binary is gone, only the
\src{x\_rect} constraint is emitted, which remains sound because the
relaxed encoding still produces a valid $z^+$.

\paragraph{Sub-LP for unknown relations.}
When \src{dep == NaN} and the upstream bounds don't immediately resolve
it (lines 152--159 handle the easy
$l_\mathrm{org} \ge u_\mathrm{pert}$ and
$l_\mathrm{pert} \ge u_\mathrm{org}$ cases), \src{encode\_dependencies}
introduces an auxiliary objective on
$\mathrm{av}[\mathrm{ind}_o + 1] - \mathrm{av}[\mathrm{ind}_p + 1]$
and calls \src{optimize!(m)} to compute a sound lower or upper bound
on the diff (lines 161--192). The Cutoff is set to 0 to terminate the
sub-LP early once a definitive sign is found. If the resulting
sub-bound is sufficiently far from zero
(\src{l\_diff > -non\_equality\_tolerance} or
\src{u\_diff < non\_equality\_tolerance}, where the default tolerance
is \texttt{1e-4}), the dependency is ratified and the corresponding
constraint pair is emitted. This sub-LP machinery is unchanged from
standard mode and is independent of the advstd techniques.

\paragraph{Interaction with the combo.}
\src{phi\_dep} is built from the \emph{same} \src{layers\_info\_dict}
that \src{relu()} populated during \src{get\_model}. Because Tech 1
already injected $N_1$'s tighter bounds and Tech 3 may have eliminated
or relaxed some neurons, the dep encoder sees the \emph{post-Tech}
view of $(u, l)$. This works correctly because:
\begin{itemize}[noitemsep]
  \item Tighter bounds make the easy resolution at lines 152--159
    fire more often, reducing sub-LP work.
  \item Eliminated binaries trigger the \src{has\_a\_o}/\src{has\_a\_p}
    guard at lines 133--134, suppressing the binary-equality
    constraint that would otherwise reference a non-existent
    \src{av[ind+2]}.
  \item Relaxed binaries (Tech 3 triangle) similarly miss
    \src{av[ind+2]} and the guard suppresses the binary constraint.
\end{itemize}

\paragraph{Caching.}
First call (\src{activation\_start == 1}) pushes each layer's resolved
\src{phi\_dep} into \src{reuse\_bounds\_conf.reusable\_deps}. On
subsequent class-pair iterations within the same Phase 2 the
encoder reuses these resolved depending-relations and skips the
sub-LP step (lines 209--212). The cache lives in
\src{help\_functions.jl}'s globals; it is cleared between Phase-2
runs.

%%─────────────────────────────────────────────────────────────────────
\section{Reproducing the safe combo end-to-end}
\label{sec:reproduce}

\subsection{Phase 1 --- solve $N_1$ once}

\begin{lstlisting}[basicstyle=\ttfamily\small,breaklines=true,frame=single]
julia run.jl --mode advanced_standard_n1 \
  --dataset mnist --model_name 3x10 \
  --model_path  /path/to/N1/model.p \
  --model_path2 /path/to/N2/model.p \
  --perturbation translation --perturbation_size "1,1" \
  --ctag 1 --ct "4,5,6" --timout 1800 \
  --n1_state_dir /path/to/n1_state \
  --use_hyper_attack       true \
  --activate_vaghgar_deps  true \
  --use_perturbed_intervals true
\end{lstlisting}

This solves $N_1$'s standard MIP (with PGD warmstart, VHAGaR-deps,
and perturbed intervals all active) and serialises
\src{n1\_layers\_info}, the primal solution, and the diff-bound
inputs into \src{n1\_state}. Repeat per $(c_s, c_t)$ pair.

\subsection{Phase 2 --- the safe combo}

\begin{lstlisting}[basicstyle=\ttfamily\small,breaklines=true,frame=single]
julia run.jl --mode advanced_standard_n2 \
  --dataset mnist --model_name 3x10 \
  --model_path  /path/to/N1/model.p \
  --model_path2 /path/to/N2/model.p \
  --perturbation translation --perturbation_size "1,1" \
  --ctag 1 --ct "4,5,6" --timout 1800 \
  --n1_state_dir /path/to/n1_state \
  --use_hyper_attack             true   \   # Cutoff + PGD Start (consumed by prev_pgd)
  --activate_vaghgar_deps        true   \   # Section 6
  --use_perturbed_intervals      true   \   # required for sound BoundTightPertRelax
  --adv_std_bound_tightening     true   \   # Tech 1, prerequisite for tau and prev_pgd
  --adv_std_zono_bounds          true   \   # zono Source A + B
  --adv_std_n2_relax_threshold   0.5    \   # Tech 3, tau = 0.5
  --adv_std_var_hint             prev_pgd   # Tech 2, prev rule + Start consensus
\end{lstlisting}

\paragraph{Healthy-stdout sanity checks.}
\begin{itemize}[noitemsep, leftmargin=2em]
  \item \src{compute\_n2\_relax\_decision!: tau=0.5, both=...
    org-only=... pert-only=... none=... (plus ... stable neurons)}
    --- non-zero \src{both + org-only + pert-only}.
  \item \src{apply\_n1\_var\_hints!: mode=PREV\_PGD, Start-consensus
    on N of M N2 binaries (pgd-silent-filled=...,
    pgd-agreed-nop=..., pgd-disagreed-withdrew=..., flipped=...)}
    --- no abort, \src{N == M}.
  \item Filename suffix on the saved CSV:
    \src{...\_N2\_advStd\_BoundTightPertRelax0.5\_zonoBounds\_varHintPrevPGD\_HyperAttackHints\_VagharDeps\_PerturbedIntervals\_cTag\{c\_tag\}}.
  \item Final \src{delta\_n2} is finite and the run terminates
    \src{OPTIMAL} or \src{TIME\_LIMIT} with a finite incumbent.
\end{itemize}

\paragraph{Unsoundness guard.}
The driver aborts before encoding if
\flag{adv\_std\_n2\_relax\_threshold}\;\texttt{>}\;0 is paired with
\flag{use\_perturbed\_intervals}\;\texttt{false}, because the
soundness argument for tiered BoundTightPertRelax requires the
inter-copy coupling to remain intact (\fileline{run.jl}{673--680}).

\end{document}
